\chapter{OBJETIVOS}

\section{Objetivos generales}
El objetivo general del trabajo consiste en ampliar el conocimiento actual sobre la aplicación de las Tensor Networks como paradigma alternativo al deep learning convencional en tareas de clasificación de imágenes, evaluando su viabilidad, rendimiento y eficiencia mediante la implementación y experimentación con distintas arquitecturas tensoriales.

Para ello, el proyecto contempla el estudio en profundidad del estado del arte en modelos de inteligencia artificial basados en Tensor Networks, la implementación de al menos dos arquitecturas tensoriales a partir de la literatura científica, su entrenamiento y evaluación utilizando recursos de computación de altas prestaciones (HPC) del servidor Lorca de la Universidad Europea, y la comparación de sus resultados frente a modelos clásicos de deep learning sobre conjuntos de datos de referencia.

De forma complementaria, el proyecto persigue adquirir experiencia práctica en el uso de herramientas computacionales propias del ámbito de la investigación en inteligencia artificial, incluyendo el manejo de entornos HPC, el uso de librerías especializadas como PyTorch y TensorKrowch, y las buenas prácticas de documentación y reproducibilidad del código científico.

\section{Objetivos específicos}

\begin{enumerate}

	\item Realizar una revisión del estado del arte de los modelos de inteligencia artificial basados en Tensor Networks aplicados a la clasificación de imágenes, analizando las arquitecturas más relevantes y sus principales resultados.

	\item Estudiar en detalle distintos modelos tensoriales para clasificación de imágenes, con el fin de comprender su funcionamiento y sus diferencias fundamentales.

	\item Implementar al menos dos modelos de inteligencia artificial basados en Tensor Networks para tareas de clasificación de imágenes, a partir de descripciones presentes en la literatura científica.

	\item Entrenar y evaluar los modelos implementados utilizando recursos de computación de altas prestaciones (HPC), familiarizándose con el uso de servidores, entornos de ejecución y herramientas básicas de gestión de trabajos.

	\item Comparar el rendimiento de los modelos basados en Tensor Networks con modelos de inteligencia artificial convencionales, empleando conjuntos de datos de referencia en clasificación de imágenes.

	\item Analizar los resultados obtenidos en términos de precisión, eficiencia computacional y comportamiento del entrenamiento, identificando las principales ventajas y limitaciones de los enfoques tensoriales.

	\item Documentar el desarrollo del proyecto, describiendo de forma clara los modelos estudiados, las implementaciones realizadas y los resultados obtenidos.

	\item Preparar y organizar un repositorio de código que recoja las implementaciones y experimentos realizados, garantizando su claridad, reproducibilidad y correcta documentación.
	
\end{enumerate}

\section{Objetivos que quedan fuera del proyecto}

\begin{itemize}

	\item Estudiar modelos de inteligencia artificial no aplicados a la clasificación de imágenes y que no estén basados en Tensor Networks.

	\item Evaluar los modelos en conjuntos de datos distintos de los considerados en el proyecto, o realizar comparaciones extensivas en múltiples benchmarks adicionales.

	\item Ejecutar los modelos en hardware cuántico real, más allá de enfoques cuántico-inspirados ejecutados en sistemas clásicos.

	\item Desarrollar aplicaciones orientadas a entornos industriales o sistemas en producción a partir de los modelos estudiados.
	
\end{itemize}

\section{Beneficios del proyecto}

Este proyecto contribuye a ampliar el conocimiento empírico sobre las Tensor Networks como herramienta de aprendizaje automático. Si bien el fundamento teórico de estas arquitecturas está bien establecido en el ámbito de la física cuántica y matemática, su comportamiento práctico en tareas de clasificación de imágenes sigue siendo un área activa de investigación. El trabajo aportará evidencia experimental concreta sobre estas cuestiones.

Desde el punto de vista computacional, el proyecto favorece el desarrollo de modelos de IA más eficientes en el uso de parámetros. Los modelos basados en Tensor Networks permiten, en principio, representar funciones complejas con un número significativamente menor de parámetros que las redes neuronales profundas convencionales. Esto tiene implicaciones directas sobre el coste energético y computacional del entrenamiento, un aspecto de creciente relevancia en la investigación en inteligencia artificial.

La realización de este proyecto supone una experiencia muy íntegra en investigación aplicada: comprende la revisión crítica de literatura científica, el diseño e implementación de experimentos reproducibles, el uso de infraestructuras HPC y el aprendizaje de herramientas y flujos de trabajo propios del ámbito profesional de la investigación en IA. Estos conocimientos y competencias son directamente transferibles a entornos académicos y laborales en el campo de la física y la computación científica.

